% Preamble - - - -
\documentclass[12pt]{article}
% Variables - - - -
\def\course{QSS 20}
\def\assignment{Literature Review}
\def\byline{Transit and the World Cup, Greater Boston MSA}
\def\me{Roshan Karim}
\ifx\byline\empty
\def\subject{\course}
\else \def\subject{\course: \byline}
\fi
% - - - - Variables
% Commands - - - -
\newcommand{\opscale}[2]{\xrightarrow{#1r_#2 \text{→ } r_#2}}
\newcommand{\opreplace}[3]{\xrightarrow{r_#1 #2r_#3 \text{→ } r_#1}}
\newcommand{\opitc}[2]{\xrightarrow{r_#1 \text{↔} r_#2}}
% - - - - Commands
% Default Packages - - - -
\usepackage[letterpaper, total={6.5in, 9in}]{geometry}
\usepackage{fancyhdr}
\setlength{\headheight}{16pt}
\usepackage{titling}
\usepackage{citation-style-language}
\cslsetup{style = american-sociological-association}
\addbibresource{references.json}
\addbibresource{references.bib}
\usepackage{xurl}
\Urlmuskip=0mu plus 1mu
\setlength{\emergencystretch}{2em}
\usepackage[pdftitle={\assignment},
    pdfsubject={\subject},
    pdfauthor={\me},
    pdfproducer={LaTeX},
    pdfcreator={LuaLaTeX},
    colorlinks,
    citecolor=black,
urlcolor=blue]{hyperref}
\usepackage{graphicx}
\graphicspath{{./images/}}
\usepackage{parskip}
\usepackage[all]{nowidow}
% - - - - Default Packages
% Unique Packages - - - -
\usepackage{amsmath}
\usepackage{txfonts}
\usepackage{fontspec}
\setmainfont{Charis SIL}
\usepackage{bm}
\usepackage{tikz}
\usepackage{pgfplots}
\pgfplotsset{compat=newest}
% - - - - Unique Packages
% Header - - - -
\pagestyle{fancy}
\lhead{\course: \assignment}
\chead{\me}
\rhead{\today}
% - - - - Header
% - - - - Preamble
\begin{document}
% Title - - - -
\begingroup
\centering
\ifx\byline\empty
\LARGE\assignment\\
\else\LARGE\assignment\\\large\byline\\[1.5em]
\fi
\endgroup
% - - - - Title
% Content - - - -
\section{Introduction}
I began this literature review considering these three research questions:
\begin{enumerate}
    \item Are Greater Boston MSA neighborhoods with greater socioeconomic disadvantage disproportionately affected by unreliable transit?
    \item How does Greater Boston MSA transit accessibility shape access to employment opportunities?
    \item Do weather-related transit disruptions in the Greater Boston MSA disproportionately affect disadvantaged communities?
\end{enumerate}
I settle on a synthesis of these three research questions to better fit this course and contribute to the literature.
My research question is as follows:
How do mega-events impact socioeconomic disparities in public transportation?
In the curated source list below, I place emphasis on the sources that best fit this synthesis.
\section{Curated Source List}
\subsection{Theory}
\begin{enumerate}
    \item \fullcite{martensTransportJusticeDesigning2016}
    \item \fullcite{lucasTransportSocialExclusion2012}
    \item \fullcite{pereiraDistributiveJusticeEquity2017}
\end{enumerate}
\subsection{Methodology}
\begin{enumerate}
    \item \fullcite{allenSizingTransportPoverty2019}
    \item \fullcite{draperDefinitionsRailwayPerformance2017}
    \item \fullcite{fuNewPerformanceIndex2007}
    \item \fullcite{jeneliusPublicTransportExperienced2018}
    \item \fullcite{lyuAnalyzingImpactService2026}
    \item \fullcite{vanoortIncorporatingServiceReliability2014}
    \item \fullcite{welchMeasureEquityPublic2013}
\end{enumerate}
\subsection{Context-Specific Research}
\begin{enumerate}
    \item \fullcite{pereiraTransportLegacyMegaevents2018}
    \item \fullcite{pereiraTransportLegacyMegaevents2020}
\end{enumerate}
\subsection{Justification}
I selected sources that fit my research question, were well-cited and supported, and were largely peer-reviewed.
From theory, I select \citet{martensTransportJusticeDesigning2016}, which provides the foundation of transport justice.
\citet{lucasTransportSocialExclusion2012} illustrates the evolution of transport-driven social exclusion over time.
Finally, \citet{pereiraDistributiveJusticeEquity2017} analyzes how transportation relates to distributive justice.

Further, I needed robust methodological resources to understand how best to conduct this research.
I looked to \citet{fuNewPerformanceIndex2007} for how to measure performance, \citet{welchMeasureEquityPublic2013} for how to measure equity, and \citet{lyuAnalyzingImpactService2026} for how to measure satisfaction.
\citet{draperDefinitionsRailwayPerformance2017} provides an overview of performance metrics, and \citet{allenSizingTransportPoverty2019} provides a quality socioeconomic status index.
\citet{jeneliusPublicTransportExperienced2018} and \citet{vanoortIncorporatingServiceReliability2014} help integrate service reliability with performance through adjusted travel time.

Finally, I took advantage of context-specific research to understand the impact of the World Cup (proxy for large scale events and/or disasters) on transit equity.
\citet{pereiraTransportLegacyMegaevents2018} and \citet{pereiraTransportLegacyMegaevents2020} provide this valuable context through study of the Rio Olympics.

Together, these sources will inform my novel transit dependability measure (a research-informed synthesis of frequency, adjusted travel time, and route/destination diversity).
I will use this measure to compare socioeconomic disparities in dependability before, during, and after the World Cup.

\section{Mapping the Intersection of Sociology and Transit}
The generated mind map shows a clear progression from theory to application.
Majorly, sources are clustered around theory and methodology, with a focus on performance metrics, socioeconomic status, and improvement techniques.
Several sources bridge disparate themes. \citet{pereiraDistributiveJusticeEquity2017} associates \citet{martensTransportJusticeDesigning2016} with applied accessibility research.
Such research (e.g. \citet{allenSizingTransportPoverty2019}) links transit accessibility to measurable social outcomes. 
\citet{pereiraTransportLegacyMegaevents2018} and \citet{pereiraTransportLegacyMegaevents2020} link accessibility research with mega-events.
Though many disparate studies provide robust operational metrics, these are not adequately linked to transit accessibility, which reduces the quality of the aforementioned studies \cite{draperDefinitionsRailwayPerformance2017,fuNewPerformanceIndex2007,lyuAnalyzingImpactService2026}.
\begin{center}
    \includegraphics[width=0.5\textwidth]{images/mindmap.png}

    \small Figure 1: NotebookLM Mind Map
\end{center}
The major gap in the literature is our understanding of the impacts of mega-events on transit equity in the short-to-medium term, specifically those caused by dependability challenges.
Work on mega-events and transportation focuses on long-term infrastructure changes, ignoring daily disruptions that impact vulnerable individuals.
In that vein, no study has been conducted on how operational decisions during mega-events contribute to disparities.
I hope to fill this gap by focusing on the short term with close study of operational decision-making and schedule modifications.
\section{Literature Review}
Discourse in transportation research has increasingly focused on the role of public transportation in facilitating access to opportunity.
Rather than measuring the efficiency, speed, or satisfaction with which transit moves its users, literature examines whether transit systems allow people to reach essential destinations.
As an example, \citet{lucasTransportSocialExclusion2012} studies how transport disadvantage inhibits participation in everyday life.
Likewise, \citet{allenSizingTransportPoverty2019} study transport poverty by examining how low- to moderate-income households face difficulty reaching employment.
In the theory of transit justice, \citet{martensTransportJusticeDesigning2016} and \citet{pereiraDistributiveJusticeEquity2017} argue that transit systems should be evaluated using the principles of distributive justice rather than efficiency metrics.
In other words, transportation systems should prioritize accessibility to underserved communities over other service goals.
These studies represent a shift of focus from infrastructure and mobility to equitable access to opportunity.

Because accessibility cannot be observed directly, researchers have developed several methods to approximate it.
\citet{allenSizingTransportPoverty2019} quantify transport poverty by measuring accessibility to employment in low- to moderate-income neighborhoods.
\citet{welchMeasureEquityPublic2013} incorporate connectivity, service frequency, capacity, and accessibility into a graph-theoretic equity index that measures how fairly service is distributed regionally.
Yet these approaches are limited by a strict focus on static transit schedules, and they lack understanding of service reliability and dependability.

Concurrently, many researchers focus exclusively on transit performance. \citet{fuNewPerformanceIndex2007} propose the Transit Service Indicator, integrating frequency, coverage, travel time, and demand.
Recently, \citet{lyuAnalyzingImpactService2026} show that service frequency and on-time performance independently influence ridership (a proxy for user satisfaction).
These studies rely primarily on computer-readable General Transit Feed Specification (GTFS) data and statistical modeling to evaluate performance, and they provide a solid assessment of the quality with which transit systems implement service schedules.
But it is clear that operational research and justice research have diverged, as these metrics, despite their vital importance to a true understanding of transit access, are not integrated with justice research by any study.

The literature on mega-events begins to associate these traditions by examining how large transit events reshape socioeconomic disparities in transit access.
\citet{pereiraTransportLegacyMegaevents2018} finds that despite Rio's transit investments for the World Cup and Olympics, accessibility was reduced for low-income groups because of concurrent service cutbacks.
This is a recurrent theme: \citet{pereiraTransportLegacyMegaevents2020} concludes that large events tend to reinforce existing inequities despite public investment.
Despite these contributions, this research is limited by its examination of one city (Rio) and its focus on infrastructure and spending, more permanent changes, as opposed to operational decisions.
In examining the 2026 FIFA World Cup in the Greater Boston metropolitan statistical area, I hope to address this gap by examining whether event-driven changes to public transit operations affect equity in transit access.
\section{Three Claims}
\begin{enumerate}
    \item \textbf{CLAIM:} \citet{lucasTransportSocialExclusion2012} defines transport-related social exclusion as a multidimensional process where transport disadvantage (e.g., no car access) compounds with social disadvantage (e.g., low income) to prevent participation in the community.
    \\ \textbf{VERDICT:} \textsc{partially accurate}
    \\ \textbf{COMMENT:} This claim is largely accurate but majorly reductive.
    \citet[106]{lucasTransportSocialExclusion2012} first defines social exclusion broadly as ``the inability to participate in normal relationships and activities,'' but also as ``the lack or denial of resources, rights, good, and services.''
    In associating social exclusion with transportation, the authors illustrate the multidimensional nature of the relationship, and at no point claim that it is as simple as transport disadvantage `compounding' with social disadvantage.
    In truth, access to transportation or the lack thereof does not have a simple relationship with social disadvantage, but one that varies based on the individuals involved.
    \item \textbf{CLAIM:} \citet{fuNewPerformanceIndex2007} propose a TSI that integrates frequency, service hours, and travel time relative to automobile travel.
    \\ \textbf{VERDICT:} \textsc{accurate}
    \\ A core component of the Transit Service Indicator is the calculated door-to-door travel time relative to the same value for automobile travel.
    This claim by NotebookLM does a good job of summarizing the TSI, though it is of course not comprehensive.
    \item \textbf{CLAIM:} \citet{allenSizingTransportPoverty2019} create a combined measure of neighborhood SES by standardizing and equally weighting four census variables: Unemployment Rate, Log of Median Household Income, the Low Income Measure, and the Low Income Cut-Off.
    \\ \textbf{VERDICT:} \textsc{accurate}
    \\ This is exactly what the authors do.
    Their neighborhood SES measure is the standardized average of the four variables, though they interpret all but income negatively.
\end{enumerate}
\section{Research Question}
Based on the above review of the literature, I propose the following research question:
How do mega-events impact socioeconomic disparities in public transportation?
I will focus on the impact of the 2026 FIFA World Cup on communities in the Greater Boston metropolitan statistical area.
I will use static and real time GTFS data to calculate transit dependability before, during, and after the World Cup.
I will apply a `differences-in-differences' approach, examining disparities in access to quality transit both pre- and post-event.
I will evaluate the socioeconomic status of communities using the methodology of \citet{allenSizingTransportPoverty2019} and publicly available data from the American Community Survey.
For transit data, I will use my archive of publicly-available GTFS-realtime data as well as the GTFS static files published by the Massachusetts Bay Transportation Authority.
This question represents a synthesis of my three me-search questions, combining study of adverse events with that of socioeconomic disadvantage.
\section{Workflow}
I enjoyed working with NotebookLM, and I made extensive use of its tooling, including the fast research function.
That tool in particular was excellent at performing cursory searches and identifying viable themes while ruling out those which had already been studied.
This helped me at the beginning when I was feeling lost on where to start, and saved me a lot of time that I would have spent mucking around and finding a path.

The fast research tool would ultimately be my own downfall, though, as I found that NotebookLM was quick to dump large swaths of sources into my notebook for manual pruning.
I spent so much time reading through a bunch of sources that had no chance of being useful to my research.
This was ultimately without consequence, as in following the assignment template, I had begun by uploading sources that I found by myself, and I ended up with a lot of really great content.
In retrospect, I should have known better than to rely on fast research for finding relevant papers, but I was initially excited by the tool and caught on to its downsides too late.
I would be able to save a lot of time by ignoring/blocking fast research tool calls beyond that initial cursory search in the future.

\textbf{NotebookLM URL:} \url{https://notebooklm.google.com/notebook/06acac37-f6fc-4bab-8488-5a98814b6db9/preview}

\printbibliography
% - - - - Content
\end{document}